\documentclass[10pt,a4paper]{article}
\usepackage[utf8]{inputenc}
\usepackage[french]{babel}
\usepackage[T1]{fontenc}
\usepackage{geometry}
\geometry{a4paper, margin=15mm}
\usepackage{booktabs}
\usepackage{xcolor}
\usepackage{hyperref}

\definecolor{primary}{RGB}{37, 99, 235}
\definecolor{dark}{RGB}{15, 23, 42}

\title{\vspace{-1.2cm}\textbf{\Large Projet 10 : Système de Recommandation Multimédia}\\
\large Bilan Complet : Pipeline DataOps, Architecture Deep Learning \& Résultats (Tâches 1 \& 2)}
\author{\textbf{Fedi BERRIMA} \and \textbf{Nadia AKHSAY}\\
\textit{Master Data Science 5 (M2) --- Module Deep Learning 2026--2027}}
\date{\vspace{-0.5cm}\today}

\begin{document}
\maketitle

\section{Bilan Chiffré de la Collecte Multimodale}
\begin{itemize}
    \item \textbf{Films analysés :} 9 742 films MovieLens croisés avec 9 621 films TMDB et 9 544 films Kaggle.
    \item \textbf{Taux de couverture multimodale :} \textbf{99.0\,\%} des films possèdent à la fois leur affiche, leur synopsis et leurs métadonnées complètes.
    \item \textbf{Volume d'images validées :} Plus de 14 000 affiches JPEG au format standard w342 vérifiées via Pillow.
    \item \textbf{Volume de notes :} 100 836 interactions utilisateurs découpées chronologiquement (77 048 Train, 9 631 Val, 9 631 Test) avec 182 films réservés exclusivement au test de démarrage à froid (\textit{Cold-Start}).
\end{itemize}

\section{Répertoire des Fichiers Python (collecte/)}
\begin{table}[h]
\centering
\small
\begin{tabular}{@{}lp{5.5cm}p{7cm}@{}}
\toprule
\textbf{Fichier Python} & \textbf{Rôle Opérationnel} & \textbf{Détails \& Sorties} \\ \midrule
\texttt{movielens.py} & Téléchargement \& validation MovieLens & Archive zip, vérification MD5, extraction \texttt{ratings.csv} et \texttt{links.csv}. \\
\texttt{tmdb.py} & Client API asynchrone (\texttt{aiohttp}) & 20 req/s, gestion des codes 429, mise en cache \texttt{movies.jsonl}. \\
\texttt{posters.py} & Téléchargement des affiches & Vérification du format RGB avec Pillow, calcul du hachage SHA-1. \\
\texttt{kaggle\_meta.py} & Secours hors-ligne (\textit{Fallback}) & Lecture de \texttt{movies\_metadata.csv} (34 Mo), comble les valeurs manquantes. \\
\texttt{normalize.py} & Harmonisation des métadonnées & Encodage des 20 genres en colonnes multi-hot (\texttt{g\_*}), normalisation langues (ISO). \\
\texttt{build.py} & Assemblage final \& Découpages & Génère \texttt{movies.csv}, \texttt{ratings.parquet} et le rapport qualité. \\
\texttt{config.py} \& \texttt{.env} & Configuration \& Clé API & Gestion des chemins et stockage sécurisé de la clé d'API TMDB. \\ \bottomrule
\end{tabular}
\end{table}

\section{Algorithmes de Deep Learning et Justification des Choix}
\begin{table}[h]
\centering
\small
\begin{tabular}{@{}lp{4cm}p{8cm}@{}}
\toprule
\textbf{Modalité} & \textbf{Algorithme Retenu} & \textbf{Pourquoi ce choix ? (Justification technique)} \\ \midrule
\textbf{Affiche (Image)} & \textbf{ResNet-50 (CNN)} & Poids gelés en phase 1. Extrait l'ambiance visuelle et le style graphique sans surapprentissage. \\
\textbf{Synopsis (Texte)} & \textbf{Sentence-BERT} (\texttt{MiniLM}) & Bi-encodeur optimisé pour la similarité cosinus. Comprend le sens réel de l'histoire. \\
\textbf{Métadonnées} & \textbf{Réseau Dense (MLP)} & Compresse les 20 genres multi-hot, la durée et la popularité en un vecteur de 64 dimensions. \\
\textbf{Couche de Fusion} & \textbf{Dense + ReLU} (Concat 320d $\to$ 128d) & Combine les 3 modalités pour former l'empreinte vectorielle unique du film. \\
\textbf{Scoring / Cold Start} & \textbf{Embedding Utilisateur + Biais} & Permet de prédire une note même pour un film sans aucune note historique. \\ \bottomrule
\end{tabular}
\end{table}

\end{document}
